\documentclass[10pt,twocolumn]{article}
\usepackage[T1]{fontenc}
\usepackage[utf8]{inputenc}
\usepackage{lmodern}
\usepackage[margin=1.8cm,top=2.4cm,bottom=2cm,columnsep=0.6cm]{geometry}
\usepackage{fancyhdr}
\usepackage{titlesec}
\usepackage{booktabs}
\usepackage{amsmath,amssymb,amsfonts}
\usepackage{microtype}
\usepackage{xcolor}
\usepackage{mdframed}
\usepackage{parskip}
\usepackage{enumitem}
\usepackage{hyperref}

\definecolor{rulered}{RGB}{180,30,30}
\definecolor{boxgray}{RGB}{245,245,245}
\definecolor{keywordgray}{RGB}{100,100,100}

\pagestyle{fancy}
\fancyhf{}
\fancyhead[L]{\textsc{\small Claude Mag}}
\fancyhead[C]{\small\textit{ICML 2026 Special Edition: Uncertainty \& Conformal Prediction}}
\fancyhead[R]{\small 2026-07-08}
\fancyfoot[C]{\small\thepage}
\renewcommand{\headrulewidth}{0.8pt}

\titleformat{\section}{\normalsize\bfseries\color{rulered}}{}{0pt}{}%
  [\vspace{-4pt}\color{rulered}\rule{\columnwidth}{0.4pt}]
\titlespacing{\section}{0pt}{8pt}{4pt}

\hypersetup{colorlinks=true,urlcolor=rulered,linkcolor=black}

\begin{document}
\setlength{\parindent}{0pt}
\setlength{\parskip}{4pt}

{\small\color{keywordgray}\textbf{Keywords:}
  conformal prediction \textbullet{}
  label efficiency \textbullet{}
  sequential stopping rules \textbullet{}
  calibration cost}\par\vspace{4pt}

{\LARGE\bfseries\raggedright
  Provably Label-Efficient Conformal Prediction}\par\vspace{4pt}

{\large\itshape\raggedright
  A Sequential Stopping Rule Cuts Calibration Labeling Cost by 41\%
  Without Sacrificing Coverage}\par\vspace{6pt}

{\small\textbf{By} Andrew Ilyas, Joonhyuk Ko, Jingwu Tang, Zhiwei Steven
  Wu \& Jiahao Zhang
  (Massachusetts Institute of Technology, Carnegie Mellon
  University)}\par\vspace{2pt}

{\small\color{keywordgray}ICML 2026 Poster}
\par\vspace{2pt}\noindent{\color{rulered}\rule{\columnwidth}{0.6pt}}\vspace{6pt}

Conformal prediction's finite-sample coverage guarantee is usually
presented as free once a calibration set exists---but collecting those
labels rarely is. This paper treats calibration labels as a costly,
queryable resource and designs an online stopping rule that
automatically balances labeling cost against prediction-set efficiency.
The rule matches the best \emph{fixed} calibration size chosen with
hindsight, cuts total labeling cost by \textbf{41.4\% $\pm$ 2.3\%}, and
provably preserves exact coverage despite deciding, on the fly, when to
stop labeling.

\begin{mdframed}[backgroundcolor=boxgray,linecolor=rulered,linewidth=1pt,
    innerleftmargin=6pt,innerrightmargin=6pt,
    innertopmargin=6pt,innerbottommargin=6pt]
\textbf{\small AT A GLANCE}\par\vspace{4pt}
\begin{description}[leftmargin=0pt,labelsep=4pt,itemsep=2pt,parsep=0pt,
    font=\small\bfseries]
  \item[Problem:] Conformal calibration-set size is treated as a free
    given, but labels needed to build it are often expensive to
    collect.
  \item[Why it matters:] Medical labeling, expert annotation, and
    expensive simulation make calibration cost the dominant cost of
    deploying conformal prediction in practice.
  \item[Method:] An online, cost-aware stopping rule for sequential
    label queries that halts once further labels are provably not
    worth their cost.
  \item[Result:] Matches the best fixed calibration size in hindsight;
    cuts total labeling cost by 41.4\% $\pm$ 2.3\% empirically.
  \item[Evidence:] Benchmark classification and regression tasks
    comparing adaptive stopping against fixed calibration-set-size
    baselines, with coverage validated throughout.
\end{description}
\end{mdframed}

\section{The Big Picture}

Split conformal prediction fixes a calibration set of size $n$ up
front; more calibration data generally tightens prediction sets, with
diminishing returns, and every additional label carries a real
acquisition cost. Conformal theory has always treated $n$ as given,
leaving practitioners to guess a calibration budget with no
statistical guidance. This paper asks the practical question directly:
given unlabeled examples arriving one at a time and the ability to
query a label at a cost, when should calibration stop and prediction
begin?

\section{Why Existing Approaches Fall Short}

Practitioners today either guess a calibration-set size in advance,
risking wasted labeling budget or an under-calibrated, oversized
prediction set, or they collect labels until budget runs out with no
statistical link between that stopping point and the tradeoff they
actually care about. Naively stopping ``when the sets look tight
enough'' is an adaptive, data-dependent rule that would ordinarily void
the exchangeability conformal prediction relies on---so ad hoc adaptive
stopping has been avoided rather than solved.

\section{Core Mechanism}

The authors formalize a total-cost objective combining the cumulative
cost of labels queried so far with the expected inefficiency
(prediction-set size) of the conformal predictor that would result from
stopping now. They design a sequential stopping rule that halts
querying once the estimated marginal benefit of one more label---in
expected set-size reduction---no longer exceeds its cost. The delicate
technical contribution is showing that this data-dependent stopping
decision can be constructed so the resulting calibration set still
yields a conformal predictor with exact finite-sample coverage,
regardless of the adaptive stopping.

\section{Technical Formulation}

With per-label cost $c$ and a weight $\lambda$ trading off cost against
expected prediction-set size $\mathbb{E}[|C(X)|]$, define the total
objective after $m$ queried labels as
\[
J(m) = m \cdot c + \lambda \cdot \mathbb{E}\big[|C_m(X)|\big],
\]
where $C_m$ is the conformal predictor calibrated on the first $m$
labels. The stopping time is
\[
\tau = \inf\{m : \widehat\Delta(m) \le 0\},
\]
where $\widehat\Delta(m)$ estimates the marginal value of querying
label $m+1$. Under mild regularity conditions, $\mathbb{E}[J(\tau)]$
matches $\min_m J(m)$---the best fixed calibration size chosen with
hindsight---up to lower-order terms, while
$P(Y \in C_\tau(X)) \ge 1-\alpha$ holds exactly despite $\tau$ being a
data-dependent stopping time.

\section{Learning or Inference Procedure}

\begin{enumerate}[itemsep=2pt,parsep=0pt]
  \item Observe an unlabeled example; decide whether to query its
    label based on the current estimate of marginal value
    $\widehat\Delta(m)$.
  \item If $\widehat\Delta(m) > 0$, query the label, add it to the
    calibration set, and increment $m$.
  \item If $\widehat\Delta(m) \le 0$, stop querying at $\tau = m$.
  \item Calibrate the conformal predictor on the $\tau$ collected
    labels and deploy it with the standard coverage guarantee.
\end{enumerate}

\section{What the Guarantee Says}

Despite $\tau$ being chosen adaptively from the data observed so far,
the resulting conformal predictor satisfies
$P(Y \in C_\tau(X)) \ge 1-\alpha$ exactly, in finite samples. Separately,
the expected total cost at the stopping time matches the best
achievable cost under a fixed calibration budget selected with
hindsight knowledge of the data.

\section{Experimental Findings}

Across benchmark classification and regression tasks, the proposed
stopping rule reduces total labeling cost by \textbf{41.4\% $\pm$ 2.3\%}
relative to fixed calibration-set-size baselines chosen without this
guidance, while maintaining nominal coverage throughout the evaluated
tasks.

\section{Ablations and Interpretation}

\textbf{Cost sensitivity:} As per-label cost $c$ increases relative to
the set-size weight $\lambda$, the stopping rule halts earlier,
trading a somewhat larger prediction set for lower total cost---exactly
the intended behavior.

\textbf{Robustness to misestimation:} Errors in estimating
$\widehat\Delta(m)$ affect the total-cost optimality of the stopping
point but never the coverage guarantee, which is preserved by
construction regardless of how $\tau$ is chosen.

\textbf{Practical takeaway:} Calibration-set size, long treated as a
free design choice in conformal prediction, can instead be solved for
directly as a cost-aware stopping problem.

\section*{Reference}

Andrew Ilyas, Joonhyuk Ko, Jingwu Tang, Zhiwei Steven Wu, Jiahao Zhang.
\textbf{Provably Label-Efficient Conformal Prediction}.
ICML 2026.
\url{https://icml.cc/virtual/2026/poster/65796}

\end{document}
